% エントリ 17(2026-11-10): m = 3(433)の超放射の印
\section{[E17, 2026-11-10] m = 3(433)の超放射の印と、433 の相対論的な表をやめたこと}\label{sec:e17}

「物理とコードは整合しているか」を確かめたとき(ユーザーの問い)、印で示していないずれが一つ見つかった。臨界が $\chi_c^{(433)}<\chi_0\le\chi_c^{(322)}$ の BH を、関数は「超放射しない」(印 0)として $\chi_0$ を返すが、実際には 433($m=3$)の超放射でゆっくり回転を失いうる。これに印をつけた(ユーザー指示「Flag it now」)。計算は \texttt{calc/e17\_m3\_flag.py}、出力は \texttt{calc/out/e17.json}。

\subsection{印のつけ方}

返すスピンが 433 の臨界より上で、433 が 1 個の量子から(自由場で)年齢のうちに飽和しうる($\gamma_{433}T>\ln[(\chi-\chi_c^{(433)})GM^2/3]$)点に印をつける。
\begin{itemize}
\item 印 0(超放射しない、$\chi_0$ を返す)は印 3(模型の外)に変える。値は $\chi_0$ のまま。
\item 印 2(自由場で 322 の臨界などに止まる)は印 4(模型の外の準位が育ちうるので値は粗い)に変える。値は $m\le2$ で計算したまま。
\end{itemize}
433 の自己相互作用(吸い込みは 766)は入れていない。自由場の飽和の時間は、433 が育つ速さの目安である。433 の臨界スピンは、211・322 と同じ臨界の表(E16)に加えた。

\subsection{433 の相対論的な表をやめた}

433 の臨界を表にしたとき、$\alpha=0.30$--$0.314$ で「どのスピンでも超放射しない」という誤った結果が出た。原因は E4 の 433 の表で、$\alpha\approx0.3$ では $\mathrm{Im}\,\omega/\mathrm{Re}\,\omega\sim10^{-16}$ と倍精度の限界にあり、高いスピンで連分数法の根が丸めの誤差に支配される(表~\ref{tab:e17prec})。表を補間すると、$\alpha=0.30$、$\chi=0.995$ で率が負になっていた。433 の率は印の判定にしか使わないので、相対論的な表をやめて解析式(1908.10370 の式 2.32)にした。解析式はこの範囲で多倍長の値の 1.5--2.3 倍で、印を多めにつける側にずれる。あわせて、臨界を探す関数が表の端(0.995)より上を外挿で調べていたのを、表の中で探すように直した。

\begin{table}[H]
\centering\small
\begin{tabular}{ccccc}
\toprule
$\alpha$、$\chi$ & 倍精度 & 多倍長 & 解析式 & $|\mathrm{Im}\,\omega/\mathrm{Re}\,\omega|$\\
\midrule
0.30、0.995 & $8.56\times10^{-17}$ & $2.28\times10^{-17}$ & $5.27\times10^{-17}$ & $7.6\times10^{-17}$\\
0.32、0.95 & $4.91\times10^{-17}$ & $4.93\times10^{-17}$ & $7.62\times10^{-17}$ & $1.5\times10^{-16}$\\
0.35、0.95 & $1.843\times10^{-16}$ & $1.834\times10^{-16}$ & $2.71\times10^{-16}$ & $5.3\times10^{-16}$\\
0.40、0.95 & $1.2638\times10^{-15}$ & $1.2634\times10^{-15}$ & $1.69\times10^{-15}$ & $3.2\times10^{-15}$\\
\bottomrule
\end{tabular}
\caption{433 の $M\,\mathrm{Im}\,\omega$(Leaver の連分数法)。}
\label{tab:e17prec}
\end{table}

\begin{result}
\begin{itemize}
\item \textbf{433 が飽和する時間}($\mu=0.71\times10^{-12}\eV$、$\chi_0$ は 433 と 322 の臨界の中点): $45\Msun$ で $3.9\times10^{10}\yr$、$50\Msun$ で $7.3\times10^9\yr$、$55\Msun$ で $1.6\times10^9\yr$、$60\Msun$ で $4.2\times10^8\yr$。6 Gyr のうちに効くのは $\sim55\Msun$ 以上。
\item \textbf{印の割合}(Beta(0.69, 2.28) の 1 世代目、6 Gyr、$f=10^{12}\GeV$): 433 のために印 3 になる標本は、$8$--$50\Msun$ で $\mu=0.71\times10^{-12}\eV$ なら 0.09\%、$0.89\times10^{-12}\eV$ なら 2.6\%。$8$--$80\Msun$ ではそれぞれ 4.2、4.6\%。重い BH や大きい $\mu$ では無視できない。
\item 単体版の書き出しの確かめに、印 0 の点で 433 が年齢のうちに飽和しないことと、印の例($50\Msun$、$\chi_0=0.485$ で印 3、$45\Msun$、0.381 で印 0)を加えた。すべて通る。
\end{itemize}
\end{result}

\begin{caveat}
433 の印は「$m=3$ の降下は模型に入っていない」ことを示すだけで、値は直していない。直すには 433 とその吸い込み(766 など)の運動論が要る。E15 の地図と増加率のうち $\alpha=0.30$ の点は、誤った 433 の表を使っていた(印の割合への影響は E16 から変わっていない: Beta の 1 世代目、$8$--$50\Msun$ で印 4 は 2.3、0.4\%)。
\end{caveat}
